\begin{center}
    \LARGE \textbf{TÓM TẮT DỰ ÁN}
\end{center}
\vspace{0.5cm}

Trong bối cảnh năng lượng tiêu thụ cho các hệ thống Sưởi, Thông gió và Điều hòa không khí (HVAC) tại các tòa nhà đang chiếm tỷ trọng ngày càng lớn, việc nghiên cứu và áp dụng các thuật toán điều khiển thông minh để tối ưu hóa năng lượng là vô cùng cấp thiết. Theo nhiều báo cáo, hệ thống HVAC tiêu thụ tới hơn 40\% tổng lượng điện năng trong một tòa nhà thương mại. Do đó, việc nâng cao hiệu suất hoạt động của hệ thống này không chỉ giúp giảm chi phí vận hành mà còn góp phần quan trọng vào mục tiêu giảm phát thải khí nhà kính và bảo vệ môi trường. Thay vì sử dụng các bộ điều khiển tĩnh truyền thống dựa trên các quy tắc cứng nhắc (Rule-Based Control), đồ án này tập trung vào việc nghiên cứu và áp dụng thuật toán Học tăng cường sâu (Deep Reinforcement Learning - DRL) để xây dựng một giải pháp điều khiển tự động thông minh và thích ứng linh hoạt cho hệ thống HVAC.

Dự án đã tái hiện và tối ưu hóa thành công mô hình điều khiển lai XGB-DQN bằng cách tích hợp tác tử (agent) Deep Q-Network (DQN) với môi trường mô phỏng gần đúng (Surrogate Model) dựa trên thuật toán XGBoost. Hệ thống sử dụng một không gian trạng thái bao gồm nhiều thông số môi trường phong phú như nhiệt độ trong nhà, nhiệt độ ngoài trời, bức xạ mặt trời, độ ẩm và các chỉ số thời tiết khác. Từ đó, mạng nơ-ron sâu được huấn luyện để ước lượng giá trị của các hành động và đưa ra quyết định điều khiển tối ưu. Điểm đột phá của hệ thống là việc ứng dụng mô hình Tiện nghi Nhiệt Thích ứng (Adaptive Thermal Comfort) theo tiêu chuẩn quốc tế ASHRAE 55 để định hình hàm phần thưởng (Reward function). Khác với cách tiếp cận thông thường luôn giữ nhiệt độ phòng ở một mức cố định, tiêu chuẩn ASHRAE 55 cho phép khoảng nhiệt độ thoải mái co giãn linh hoạt dựa trên sự thay đổi của thời tiết bên ngoài và khả năng thích nghi sinh lý tự nhiên của cơ thể con người. Cách tiếp cận này giúp tác tử AI nhận được tín hiệu thưởng chính xác hơn, cân bằng hoàn hảo giữa hai mục tiêu cốt lõi là tiết kiệm năng lượng và đảm bảo sự thoải mái tối đa cho người sử dụng.

Kết quả thực nghiệm trên tập dữ liệu thời tiết thực tế xuyên suốt cả năm cho thấy những thành tựu khả quan. Tác tử DQN sau quá trình huấn luyện đã tự động học được chiến lược điều khiển khôn khéo: hệ thống biết tận dụng thời điểm thời tiết mát mẻ để mở cửa sổ đón gió tự nhiên thay vì chạy điều hòa, đồng thời linh hoạt điều chỉnh công suất làm lạnh một cách tối ưu khi nhiệt độ ngoài trời tăng cao. Cụ thể, hệ thống đạt được tỷ lệ vi phạm mức độ tiện nghi nhiệt ở mức 0\% (vượt trội hoàn toàn so với sự không ổn định của các phương pháp điều khiển cơ sở). Cùng với đó, phương pháp điều khiển bằng AI này đã chứng minh hiệu quả giảm thiểu điện năng tiêu thụ một cách đáng kể. 

Tóm lại, dự án đã khẳng định tính khả thi và hiệu quả của việc kết hợp Deep Reinforcement Learning với các tiêu chuẩn tiện nghi nhiệt hiện đại trong bài toán quản lý năng lượng. Thành công của đồ án không chỉ giải quyết triệt để bài toán điều khiển HVAC mà còn cung cấp một khuôn khổ thực tiễn, đóng góp một giải pháp hứa hẹn cho việc phát triển các nền tảng quản trị tòa nhà thông minh thế hệ mới, hướng tới sự phát triển xanh và bền vững.

\cleardoublepage